\scanpage{189}
\mspage{90}

On a ceci:
\[
(304)\qquad D_{\Delta_r}\xrightarrow{\sigma}D_{\Delta_0}
\]
application simpliciale correspondant au ``premier sommet'' d'un simplexe \unsure{ordonné}, ou à la ``source'' d'une suite de flèches
\diagram{\begin{tikzcd}[ampersand replacement=\&,column sep=small]
\bullet \arrow[r,dash] \& \bullet \arrow[r] \& \bullet \arrow[r] \& \cdots \arrow[r] \& \bullet
\end{tikzcd}}

\margin{\uline{NB} $D^*_{\Delta_r}$ est l'image inverse de $D^*_{\Delta_0}$}

\unsure{est} une application finie étale \uline{au dessus} de $D^*_{\Delta_0}$ (\uline{NB} Dans le cas ``croise-ments normaux, c'est $=$ un \unsure{rev.} étale fini sur $D_{\Delta_0}$ tt entier --- mais c'est là un fait très spécial, de $=$ que le fait que les \deleted{\unsure{diff.}} strates \unsure{soient} des variétés topologiques --- ce qui implique \unsure{donc} qu'il en est \unsure{donc|bien} de $=$ des $D_{\Delta_{\unsure{*}}}$.)

Egalement le fait que \deleted{\unsure{la}} la structure ordonnée \inserted{locale} de $D^*_{\Delta_{\unsure{*|1}}}$ sur $D^*_{\Delta_0}$ (\deleted{flèches} structure \inserted{de l'ens.} des strates locales contenant une strate locale donnée) soit \deleted{\unsure{une|celle d'une}} \inserted{isomorphe} \deleted{structure de simplexe}

\margin{188}
